\documentclass[../main.tex]{subfiles}
\begin{document}

\section{Conclusiones}

El objetivo principal de este trabajo ha sido identificar nuevas mezclas de fluidos refrigerantes que permitan sustituir al R290 mejorando su eficiencia energética en bombas de calor. Para ello, se llevó a cabo una simulación teórica con diversas mezclas compuestas por refrigerantes naturales, orientada a generar soluciones a largo plazo conformes con la normativa europea F-gas. De esta fase se obtuvo la mezcla más prometedora R1270/RE170 en proporciones másicas de 7\%/93\%, que fue posteriormente objeto de estudio y evaluación experimental.

La etapa experimental empezó con el estudio de fluidos puros R290 y RE170, posteriormente se estudió la mezcla de R290/RE170 (85\%/15\% másico) sugerida por Mikoleit et al. \cite{Mickoleit_2026} pero los resultados que se obtuvieron no presentaron mejoras significativas por lo que se descartó esta mezcla como sustituta del R290.

Posteriormente se ensayó la mezcla que resultó más prometedora de la simulación teórica, esta sí que presentó resultados prometedores ya que se observó que, al incrementar la proporción de R1270 en la matriz de RE170, el COP se mantenía igual que el del RE170 puro mientras que el VHC aumentaba. Se fue capaz de estudiar hasta las proporciones másicas 20\%/80\% (R1270/RE170), punto en el que se terminó el tiempo de ensayo aunque el COP siguiera siendo constante y el VHC siguiera aumentando. El aumento del COP y los buenos resultados de VHC sugieren que esta mezcla es una candidata viable para sustituir al R290, ya que es más eficiente y presenta una implementación relativamente sencilla, ya que la única modificación que necesita la instalación es la reconfiguración de la válvula electrónica y el cambio de refrigerante. Además, al combinar un GWP reducido con una menor demanda eléctrica y, por consiguiente, menores emisiones indirectas de [CO$_\text{2}$]$_\text{eq}$ se cumplen al completo los requisitos de la normativa F-gas.

Desde una perspectiva económica, la adopción de la composición R1270/RE170 (20\%/80\%) en sustitución del R290 en equipos comerciales de bomba de calor representaría, para una empresa fabricante, una tasa interna de retorno comprendida entre el 48\% y el 151\%, lo que indica que se trata de una inversión rentable.

En cuanto a las líneas de investigación futuras, sería conveniente ampliar el estudio a otras composiciones de R1270/RE170 para evaluar su comportamiento bajo condiciones de igual capacidad calorífica volumétrica y analizar su eficiencia energética en dichas condiciones. Asimismo, queda pendiente el estudio de la mezcla ternaria R1270/RE170/R600a en proporciones másicas 19\%/65\%/16\%, que en la fase teórica mostró un incremento notable en la eficiencia energética con una reducción moderada del VHC.


\newpage

\end{document}